% Generated from ../chapters/17-recovery-remodeling-and-regeneration.md
\chapter{Recovery, Remodeling, and Regeneration}

Several distinct processes are often compressed into the word ``healing.''

\textbf{Functional compensation} uses remaining capacity differently.

\textbf{Recovery of regulation} removes inhibition or maladaptive control.

\textbf{Remodeling} changes existing tissue architecture.

\textbf{Repair} restores continuity, often with scar or imperfect replacement.

\textbf{Regeneration} reconstructs lost tissue with appropriate cells, architecture, and integration.

These processes overlap but are not interchangeable.

\section{Why the distinction matters}

Muscle guarding can be released without creating new muscle cells. Bone can remodel under altered loading without returning to a childhood shape. A surviving neuron may regrow an axon without replacement of dead neurons. A new cell may be produced yet fail to integrate into a functional circuit.

Claims should identify the level of restoration.

\section{Adult plasticity is not binary}

Tissues do not divide simply into ``regenerative'' and ``non-regenerative.'' Adult plasticity may involve:

\begin{itemize}
\tightlist
\item
  active stem-cell niches;
\item
  quiescent cells;
\item
  mature-cell proliferation;
\item
  dedifferentiation;
\item
  transdifferentiation;
\item
  axonal sprouting;
\item
  extracellular-matrix remodeling;
\item
  immune-mediated clearance;
\item
  changes in cell state without cell replacement.
\end{itemize}

Research on partial reprogramming, mechanotransduction, and regenerative signalling demonstrates that mature tissues retain more state flexibility than once assumed. It does not demonstrate a universal latent program accessible through lifestyle.

\section{The eye as a boundary case}

The optic nerve illustrates why tissue-specific analysis is essential. If retinal ganglion cells survive but their axons are damaged, regrowth is conceptually different from replacing dead cells. Mouse experiments using targeted OSK expression have promoted axonal regeneration and restored some visual function. That finding shows that mature neuronal state is not absolutely fixed.

It does not show that meditation, fasting, or systemic relaxation can produce the same molecular intervention in humans. Nor does it solve long-distance guidance, remyelination, target selection, and circuit integration after severe loss.

The framework should remain open without becoming vague:

\begin{quote}
The absence of spontaneous recovery is not proof of metaphysical impossibility; evidence of plasticity elsewhere is not proof of recoverability here.
\end{quote}

\section{Restoration as organization plus material}

A functioning structure requires both:

\begin{itemize}
\tightlist
\item
  suitable material---cells, matrix, blood supply, and energy;
\item
  suitable organization---position, identity, connectivity, timing, and regulation.
\end{itemize}

The river metaphor emphasizes organization, but a dry riverbed still needs water. A complete longevity model must account for both.
